\documentclass[11pt,a4paper]{article}

\usepackage{fontspec}
\usepackage[margin=20mm]{geometry}
\usepackage{booktabs,tabularx,array}
\usepackage{graphicx}
\usepackage{xcolor}
\usepackage{hyperref}
\usepackage{enumitem}

\defaultfontfeatures{Ligatures=TeX,Scale=MatchLowercase}
\IfFontExistsTF{Noto Sans CJK KR}{
  \setmainfont{Noto Sans CJK KR}[
    ItalicFeatures={FakeSlant=0.18},
    BoldItalicFeatures={FakeSlant=0.18}
  ]
  \setsansfont{Noto Sans CJK KR}[
    ItalicFeatures={FakeSlant=0.18},
    BoldItalicFeatures={FakeSlant=0.18}
  ]
}{
  \IfFontExistsTF{Apple SD Gothic Neo}{
    \setmainfont{Apple SD Gothic Neo}[
      ItalicFeatures={FakeSlant=0.18},
      BoldItalicFeatures={FakeSlant=0.18}
    ]
    \setsansfont{Apple SD Gothic Neo}[
      ItalicFeatures={FakeSlant=0.18},
      BoldItalicFeatures={FakeSlant=0.18}
    ]
  }{
    \setmainfont{NanumGothic}[
      ItalicFeatures={FakeSlant=0.18},
      BoldItalicFeatures={FakeSlant=0.18}
    ]
    \setsansfont{NanumGothic}[
      ItalicFeatures={FakeSlant=0.18},
      BoldItalicFeatures={FakeSlant=0.18}
    ]
  }
}
\IfFontExistsTF{DejaVu Sans Mono}{
  \setmonofont{DejaVu Sans Mono}
}{
  \IfFontExistsTF{Menlo}{
    \setmonofont{Menlo}
  }{
    \setmonofont{Liberation Mono}
  }
}
\XeTeXlinebreaklocale "ko"
\XeTeXlinebreakskip = 0pt plus 1pt

\hypersetup{
  colorlinks=true,
  linkcolor=blue!50!black,
  urlcolor=blue!50!black,
  pdftitle={UUV MuJoCo v2.2 current content notes},
  pdfauthor={OpenAI Codex}
}

\setlength{\parskip}{0.5em}
\setlength{\parindent}{0pt}
\renewcommand{\arraystretch}{1.2}
\newcolumntype{Y}{>{\raggedright\arraybackslash}X}
\newcommand{\code}[1]{\nolinkurl{#1}}

\title{UUV MuJoCo v2.2 current 내용 정리본\\
\large 시각 자료집과 분리한 설명 중심 PDF}
\author{}
\date{2026-04-07}

\begin{document}
\maketitle

\section{문서 목적}

이 문서는 그림 위주의 자료집과 분리해서, \textbf{current 모델이 어떤 구조로 동작하는지와 각 그래프를 어떻게 읽어야 하는지}를 짧게 정리한 설명본이다. 기준 입력은 아래 세 가지다.

\begin{itemize}[leftmargin=2em]
  \item \code{uuv\_mujoco/v2.2/config/sim\_profiles.json}
  \item \code{uuv\_mujoco/v2.2/scenes/tank\_current\_scene.xml}
  \item \code{document/docsource/uuv\_v22\_current\_focus\_metrics.json}
\end{itemize}

\section{current 모델 구조}

current는 \textbf{MuJoCo built-in ellipsoid fluid} 와 \textbf{Python custom thrust / hydrostatics} 를 결합한 구조다.

\begin{itemize}[leftmargin=2em]
  \item MuJoCo built-in은 ellipsoid fluid drag, viscous fluid force / torque, \code{mj\_step()} 기반 강체 운동 적분을 맡는다.
  \item Python custom은 thrust shaping과 mixing, distributed buoyancy, restoring torque, runtime mass / CoM / inertia setup을 맡는다.
  \item thrust는 actuator만으로는 실제 추진기 특성을 충분히 표현하기 어렵고, buoyancy는 current가 원하는 distributed hydrostatics를 built-in만으로 구현하기 어렵기 때문에 Python custom으로 처리한다.
\end{itemize}

\section{thrust model 설명}

current의 thrust path는 다섯 단계로 읽으면 된다.

\begin{enumerate}[leftmargin=2em]
  \item 입력은 \([-1, 1]\) 범위의 정규화된 명령값이다. 아직 실제 추력[N]은 아니다.
  \item 1차 지연이 적용된다. \code{tau\_up}=0.04\,s, \code{tau\_down}=0.06\,s 로 rise와 decay를 다르게 둔다.
  \item deadzone과 command limit로 작은 입력과 과도한 입력을 정리한다.
  \item T200 16V 성능표를 참고한 force curve를 바탕으로 normalized input을 실제 thrust[N]로 바꾼다.
  \item 마지막으로 수직축과 yaw축에만 별도 응답 보정을 적용한다. 수직은 부력 영향이 커서, yaw는 회전 응답이 실제보다 약하게 나오기 쉬워서 따로 보정한다.
\end{enumerate}

현재 스냅샷의 핵심 값은 \code{thruster\_force\_max}=32.0\,N, \code{vertical\_thruster\_gain\_scale}=1.45, \code{yaw\_torque\_scale}=1.50, \code{deadzone}=0.002, \code{gain\_scale\_all}=2.9 이다.

\section{buoyancy model 설명}

current는 한 점에 전체 부력을 몰아주는 single-point center 대신, \textbf{4개의 buoyancy point}에 부력을 분산 적용한다.

\begin{itemize}[leftmargin=2em]
  \item \code{body\_components} 는 주로 질량 분포를 정의하는 블록이다. 여기서 CoM과 runtime inertia를 계산한다.
  \item 실제 current 실행에서 부력은 \code{buoyancy\_points} 4개에 걸린다.
  \item 핵심은 부력 총량보다 \textbf{부력이 어디에 나눠서 걸리느냐}가 자세 응답을 크게 바꾼다는 점이다.
  \item single-point center는 pitch가 과도하게 커지고 자세가 오래 흔들리기 쉽고, 4-point는 전진 시 nose-up과 release 이후 진동을 줄이기 위한 구조다.
\end{itemize}

최근 current 스냅샷 값은 \code{buoyancy\_scale}=1.002, \code{cob\_x\_offset}=0.009\,m, \code{cob\_z\_offset}=0.010\,m 이다.

\section{actual reference와 그래프 해석}

여기서 actual reference는 실제 로봇 rosbag의 센서 / 추정 응답에서 뽑은 요약값이다. 즉 실제 힘을 비교한 것이 아니라, \textbf{실제 응답 envelope} 를 비교한 것이다.

\begin{center}
\small
\begin{tabularx}{\linewidth}{>{\raggedright\arraybackslash}p{0.46\linewidth}Y}
\toprule
latest actual reference 항목 & 값 \\
\midrule
Forward peak surge [m/s] & 0.557 \\
Forward peak $|$pitch$|$ [deg] & 13.652 \\
Forward peak $|$roll$|$ [deg] & 0.836 \\
Heave depth delta [mm] & 42.339 \\
\bottomrule
\end{tabularx}
\end{center}

그래프 해석은 이렇게 읽으면 된다.

\begin{itemize}[leftmargin=2em]
  \item \textbf{peak surge}: 전진 명령에서 순간 최대 전진 속도
  \item \textbf{peak pitch / roll}: 전진 명령이 자세로 얼마나 강하게 연결되는지
  \item \textbf{depth delta}: heave 입력에서 깊이가 얼마나 변했는지
  \item \textbf{settling time}: 명령 해제 후 자세가 다시 안정되는 데 걸리는 시간
  \item \textbf{gap index}: actual reference와의 상대 오차를 한 값으로 묶은 요약 지표
\end{itemize}

\section{현재 snapshot 비교값}

아래 표는 \textbf{single-point center}, \textbf{current 4-point}, \textbf{latest actual reference} 를 같은 지표로 나란히 놓은 값이다. 이 수치는 최근 scene / proxy 수정이 반영된 current snapshot 기준이다.

\begin{center}
\small
\begin{tabularx}{\linewidth}{>{\raggedright\arraybackslash}p{0.34\linewidth}YYY}
\toprule
항목 & single-point center & current 4-point & latest actual ref \\
\midrule
Forward peak surge [m/s] & 0.586 & 0.647 & 0.557 \\
Forward peak $|$pitch$|$ [deg] & 67.751 & 20.119 & 13.652 \\
Forward peak $|$roll$|$ [deg] & 0.339 & 2.329 & 0.836 \\
Heave depth delta [mm] & 133.399 & 166.198 & 42.339 \\
Attitude release settling time [s] & 5.998 & 0.680 & -- \\
Actual gap index [\%] & 144.215 & 133.637 & -- \\
\bottomrule
\end{tabularx}
\end{center}

이 표는 single-point가 여전히 pitch와 자세 복원에서 불리하고, 최근 current 4-point snapshot도 heave와 roll 쪽은 추가 재튜닝이 필요하다는 점을 보여준다. 즉 4-point 구조 자체는 유지하되, fluid proxy와 hydrostatic 파라미터는 아직 조정 여지가 남아 있다.

\section{정리}

current 모델을 짧게 정의하면 다음과 같다.

\begin{quote}
MuJoCo는 ellipsoid fluid drag와 적분을 맡고, Python은 thrust 모델과 distributed buoyancy를 맡는다.
\end{quote}

또한 최근 수정의 방향은 다음 두 줄로 요약할 수 있다.

\begin{itemize}[leftmargin=2em]
  \item thrust 쪽은 실제 추진기 특성과 축별 응답 차이를 보정하는 방향
  \item buoyancy / fluid proxy 쪽은 자세 응답과 heave 감쇠를 실제에 가깝게 맞추는 방향
\end{itemize}

\end{document}
